% ================================================================
% ch02_dariya_sagar.tex — अध्याय २: दरिया सागर
% ================================================================

\cleardoublepage
\section*{अध्याय २: दरिया सागर}
\addcontentsline{toc}{section}{अध्याय २: दरिया सागर}

\subsection*{२.१ मंगलाचरण आ आत्म-बोध}
\addcontentsline{toc}{subsection}{२.१ मंगलाचरण आ आत्म-बोध}

\textbf{मूल भोजपुरी रचना:}

\begin{quote}
\addfontfeatures{Color=3D3428}
आदि नाम सत्य पुरुष हमारा, ताकर महिमा अगम अपारा।\\
दरिया गावै सहज सुहाई, सतगुरु कृपा मिलल अधिकाई॥
\end{quote}

\textbf{भोजपुरी भावार्थ:} हमर मूल आधार सत्य पुरुष (ईश्वर) के नाम ह, जिनकर महिमा अपार आ अगम बा। सतगुरु के कृपा से दरिया साहब सहज भाव से ई महिमा गावत बाड़ें।

\begin{sloppypar}
{\engfont \textit{English Translation:} The Eternal Name of the True One is our foundation, boundless and beyond grasp.
By the grace of the True Master, Dariya sings this praise with natural devotion.}
\end{sloppypar}

\begin{center}
{\color{bordercolor}\rule{0.4\textwidth}{0.4pt}}
\end{center}

\subsection*{२.२ संसार अनित्यता आ विवेक}
\addcontentsline{toc}{subsection}{२.२ संसार अनित्यता आ विवेक}

\textbf{मूल भोजपुरी रचना:}

\begin{quote}
\addfontfeatures{Color=3D3428}
ई जग देखत सपना जैसा, छिन में बिनसे पानी जैसा।\\
सत्य नाम संबल कर लीजै, जनम जनम के संसय छीजै॥
\end{quote}

\textbf{भोजपुरी भावार्थ:} ई संसार सपना लेखा अनित्य बा, पानी के बुलबुला लेखा छिन भर में बिला जाला। एह से सत्य नाम के सहारा ला, ताकि जनम-जनम के डर आ भ्रम खतम हो जाव।

\begin{sloppypar}
{\engfont \textit{English Translation:} This physical world is transitory like a dream, disappearing in a moment like water.
Hold fast to the True Name, so that the doubts of lifetimes may be dissolved.}
\end{sloppypar}

\begin{center}
{\color{bordercolor}\rule{0.4\textwidth}{0.4pt}}
\end{center}

\subsection*{२.३ अमर लोक आ सत्य पुरुष}
\addcontentsline{toc}{subsection}{२.३ अमर लोक आ सत्य पुरुष}

\textbf{मूल भोजपुरी रचना:}

\begin{quote}
\addfontfeatures{Color=3D3428}
अमर लोक में पुरुष बिराजे, ताकर तेज अमित छबि छाजै।\\
नहिं वहाँ धूप न छाया भाई, सहज रूप सब रहल समाई॥
\end{quote}

\textbf{भोजपुरी भावार्थ:} अमर लोक में सर्वव्यापी परमात्मा विराजमान बाड़ें, जिनकर प्रकाश आ तेज अमित बा। ओहिजा कवनो धूप या छाँव के द्वंद्व नइखे, सब कुछ सहज आनंद में समाइल बा।

\begin{sloppypar}
{\engfont \textit{English Translation:} In the Eternal Realm abides the Supreme Spirit, glowing with boundless splendor.
There is neither heat nor shadow there; all creation rests in spontaneous harmony.}
\end{sloppypar}

\begin{center}
{\color{bordercolor}\rule{0.4\textwidth}{0.4pt}}
\end{center}

\subsection*{२.४ संसार उत्पत्ति आ जीव के दशा}
\addcontentsline{toc}{subsection}{२.४ संसार उत्पत्ति आ जीव के दशा}

\textbf{मूल भोजपुरी रचना:}

\begin{quote}
\addfontfeatures{Color=3D3428}
आदि कला से भइल पसारा, पाँच तत्व आ गुन विस्तारा।\\
जीव भूलि आपन निज धामा, भरमि परल माया के धामा॥
\end{quote}

\textbf{भोजपुरी भावार्थ:} परमात्मा के आदि शक्ति से एह संसार आ पाँच तत्वन के फैलाव भइल। बाकिर जीव आपन असली घर के भुला के माया के अन्हरिया में भटक गइल।

\begin{sloppypar}
{\engfont \textit{English Translation:} From the primordial cosmic power arose the expanse of the five elements.
Forgetful of its true origin, the soul wanders amidst the maze of illusion.}
\end{sloppypar}

\begin{center}
{\color{bordercolor}\rule{0.4\textwidth}{0.4pt}}
\end{center}

\subsection*{२.५ सतनाम महिमा}
\addcontentsline{toc}{subsection}{२.५ सतनाम महिमा}

\textbf{मूल भोजपुरी रचना:}

\begin{quote}
\addfontfeatures{Color=3D3428}
सतनाम सम और न दूजा, छाड़ि कपट कर सहजहिं पूजा।\\
सब साधन में नाम प्रधाणा, जाके सुमिरे कटै बंधाना॥
\end{quote}

\textbf{भोजपुरी भावार्थ:} सत्य नाम के समान दूसर कवनो साधना नइखे। मन के सब कपट छोड़ के सहज भाव से नाम सुमिरन करा। नाम ही सब साधन में श्रेष्ठ ह, जेकरा सुमिरन से सब बंधन कट जाला।

\begin{sloppypar}
{\engfont \textit{English Translation:} There is no refuge higher than the True Name; discard pretense and worship with simplicity.
Among all spiritual practices, Remembrance of Name is supreme, severing all bonds of ignorance.}
\end{sloppypar}

\begin{center}
{\color{bordercolor}\rule{0.4\textwidth}{0.4pt}}
\end{center}

\subsection*{२.६ सहज सुमिरन आ मन निरोध}
\addcontentsline{toc}{subsection}{२.६ सहज सुमिरन आ मन निरोध}

\textbf{मूल भोजपुरी रचना:}

\begin{quote}
\addfontfeatures{Color=3D3428}
सुमिरन ऐसा कीजिए, दूजा मन नहिं आय।\\
रोम रोम में नाम रमै, सुरति निरति ठहराय॥
\end{quote}

\textbf{भोजपुरी भावार्थ:} नाम सुमिरन अइसन होखे के चाहीं कि मन में कवनो दूसर विचार मत आवे। रोम-रोम में नाम समा जाव आ ध्यान पूरी तरह स्थिर हो जाव।

\begin{sloppypar}
{\engfont \textit{English Translation:} Remember the Divine in such depth that no secondary thought enters the mind.
May the Holy Name permeate every pore, stilling consciousness into unwavering focus.}
\end{sloppypar}
